% =============================================================
%  HeInSulfides_Main_Body.tex
%  Shared, class-agnostic body of Chapter 3 (Introduction through
%  Conclusions -- everything after the abstract). \input{} by
%  Chapter3/chapter3.tex. The abstract text lives separately in
%  HeInSulfides_abstract.tex, same convention as Chapters 1 and 2 --
%  see the comment there for why.
%
%  Citations use \citep/\citet against the master references.bib
%  (dissertation_preamble.tex), same as Chapters 1 and 2.
% =============================================================

\section{Introduction}\label{sec:ch3-introduction}

Sulfide minerals are nearly ubiquitous in the formation histories of economic ore deposits, 
and indeed are themselves the primary repositories for many base and precious
metals in the crust. The question of when and how they formed is
central to understanding ore-forming processes, crustal fluid
circulation, and the tectonic contexts in which economic deposits
develop. Despite decades of study, the timing of sulfide mineralization
remains contested or poorly resolved across a wide range of deposit
types, from shallow epithermal systems to deep, high temperature
magmatic deposits in mafic intrusions. This uncertainty reflects both
the genuine complexity of ore-forming systems --- which are frequently
multistage, hydrothermal, and overprinted --- and the limitations of the
geochronological tools available for dating sulfide phases directly.

\subsection{An Alternative and Compliment to Re-Os Dates in Sulfide
Geochronology}

The rhenium-osmium isotope system is the most widely applied method for
direct radiometric dating of sulfide minerals, exploiting the relatively
high Re concentrations that pyrite, molybdenite, arsenopyrite, and
pyrrhotite can accommodate. Re-Os geochronology has produced
geologically meaningful ages across a wide range of sulfide-bearing
deposit types, and remains indispensable for directly dating ore-forming
events \citep{chen2015,hnatyshin2026,schoenberg1999,ostendorf2019,selby2002}. Nevertheless, the
Re-Os system is vulnerable to a set of well-documented failure modes,
and a growing number of published ages exhibit scatter, implausible
geologic interpretations, or internal inconsistencies that motivate the
development of independent dating methods \citep{stein2003,suzuki2001,li2025,lyu2020}.

More practically, the highly specialized and costly n-TIMS method for
Re-Os dating restricts its application to a select few labs, limiting
the scale of its application \citep{ackerman2025}. Furthermore, Re-Os dating
of sulfides (other than the rare mineral molybdenite) requires
gram-quantities \citep{qi2010} and is therefore poorly suited to
resolve distinct mineralization episodes in multistage systems.

For multistage systems in particular, the co-application of independent
radiochronometers is often useful to distinguish the timings for
hydrothermal alteration from uplift, cooling, exhumation, and
potentially earlier emplacement of host rocks. For example, \citet{selby2002}
demonstrate that Re-Os molybdenite is more likely to accurately
record the age of sulfide mineralization in the Tintina Gold Belt,
Alaska, relative to the comparatively thermally-sensitive Ar-Ar mica
ages. However, molybdenite is far from a ubiquitous phase in
sulfide-bearing ores, unlike pyrite. A pyrite U-Th/He geochronometer
would expand such opportunities to co-apply radiochronometers that
interrogate processes operating at different times and temperatures to a
much more diverse and expansive array of locales that experienced
complex histories. Similar geochronometers developed for other sulfide
phases such as arsenopyrite, galena, and sphalerite may offer the
opportunity to disentangle multiple episodes of sulfide mineralization
and alteration within an ore body, allowing geologists to better
understand the sequence of events that lead to economic ore formation.

One such example of an opportunity for applying U-Th/He sulfide
geochronology is in the world\textquotesingle s largest known
platinum-group element deposit, the Bushveld Complex of South Africa,
which hosts the sulfide-bearing Merensky Reef and Platreef PGE deposits in the sulfide-bearing
upper critical zone, which for which geochronometers yield a surprising diversity of ages:
zircon U/Ph ages vary from $\approx$2057-2060 Ma, while rutile ages are significantly younger at $\approx$2053 Ma 
\citep{scoates2021}, yet Re-Os systematics show evidence for
multiple sulfide mineralization events within the same horizon
\citep{schoenberg1999}. The Re/Os and Rb/Sr isochron ages they report
are unable to resolve these distinct mineralization events from the
zircon U/Pb emplacement age. Unlike Re/Os dating, the U-Th/He system
produces distinct ages for every dated aliquot, potentially offering the
opportunity to resolve the timings of prolonged volatile activity or
late hydrothermal Re-Os disturbance.

Similarly, at the Norilsk-Talnakh deposits of the Siberian Traps, the
ore-bearing intrusions can be correlated with a specific transition in
flood basalt source and eruptive style \citep{yao2021}, yet only
three of many differentiated intrusions in the region are economic, and
the precise timing and mechanism of sulfide saturation relative to trap
emplacement is not fully understood. In the 1.1 Ga Sonju Lake Intrusion
of Minnesota, USA, stable isotopic evidence suggests that
meso/epithermal fluids reacted with primary magmatic sulfides raising
questions about when redistribution occurred and how much metal was
moved \citep{park2004}. These are deposits where independent
thermochronometric constraints capable of distinguishing primary
crystallization from later hydrothermal reworking are needed.

\subsection{The Single-Jump Migration Model and its Application to Pyrite}

The Russian school of noble gas geochronology spearheaded by
Shukolyukov and Yakubovich, approaches the problem of helium migration
in minerals through a theoretical framework that differs fundamentally
from the volume diffusion model used for to describe noble gas migration
dominant in the Western thermochronological literature \citep{dodson1973,
harrison1985,fortiergiletti1989,lovera1991,wolf1996,farley2000,lee1997}. In the Shukolyukov
"single-jump model," helium atoms occupy discrete energy minima within
the crystal structure, and migration proceeds through thermally
activated jumps between these sites \citep{shukolyukov2009,yakubovich2010}.

Within this framework, \citet{shukolyukov2012} developed the \textsuperscript{190}Pt-\textsuperscript{4}He
system as a new geochronometer for native platinum minerals,
demonstrating through thermal desorption analysis that radiogenic \He\ in
isoferroplatinum is stable to temperatures approaching the melting point
of platinum, and obtaining \textsuperscript{190}Pt-\textsuperscript{4}He ages from the Galmoenan and Kondyor
massifs consistent with independent Sm-Nd and Rb-Sr estimates. He
storage in native platinum was attributed to accumulation in vesicles
and clusters at dislocation cores and grain boundaries. Western scientists
similarly model helium migration in metals with single-jump mechanistic
single-jumps, although these single jumps are not typically modeled to
instantaneously liberate helium, instead the modeled single-jump allows
the atom to spontaneously travel a distance of some angstroms with some
frequency, introducing a spatial scale from which a diffusivity can be
evaluated \citep{ratzke1992,costhickman1975,wilson1973}.

\citet{yakubovich2015} extends this method to sperrylite (PtAs\textsubscript{2}) a
rare mineral found in magmatic PGE deposits that is isostructural with
pyrite. \textsuperscript{190}Pt-\textsuperscript{4}He ages from sperrylite at the Kondyor and
Norilsk-Oktyabrsky deposits were geologically consistent, and the high
He retentivity demonstrated in this arsenide prophesies similar behavior
in pyrite.

Without providing any sulfur buffer, \citet{yakubovich2019}
interpreted ramped-heating experiments for a variety of silica glass
encapsulated sulfides under the single-jump step-heating method. Pyrite,
in particular, was found to be anomalously high retentivity of helium. 
The first application of U-Th/He pyrite geochronometry to a natural ore
deposit was reported by \citet{yakubovich2019,yakubovich2020} for the Uzelga
Cu-Zn VMS deposit of the South Urals, yielding a consistent weighted
mean age of 382 ± 12 Ma that is older than regional prehnite-pumpellyite
metamorphism, demonstrating that pyrite can survive a metamorphic
overprint without complete isotopic resetting. While these results
establish the viability of pyrite as a U-Th/He geochronometer, their
analytical approach employs a simple, low-resolution monotonic heating schedule without
grain-size scaling experiments, sulfur buffers, and proton-implanted \textsuperscript{3}He
domain mapping, leaves the diffusion/migration physics of sulfides
undercharacterized. Unsurprisingly, given the rapid heating rates used in their experiments, 
their results for pyrite generally show abrupt helium loss at temperatures corresponding to the upper stability 
field of pyrite around 700°C \citet{toulminbarton1964}, and therefore are unlikely to contain substantial information about helium loss 
in natural pyrite over geologic timescales. Extending these initial observations into a robust, predictive thermochronometric framework requires the more
rigorous experimental and theoretical treatment that is one of the aims of this work. 

\subsection{Geologic Setting}
\label{sec:ch3-geologic-setting}

Pyrite and Arsenopyrite for the geochronologic study was collected from two deposits that strongly 
differ in the timing and conditions of their formation: the Archean Kidd Creek volcanogenic massive sulfide (VMS)
deposit, Ontario, Canada, and the Permo-Triassic orogenic gold system at
Gympie, Queensland, Australia. Pairing a $\sim$2.7 Ga Cu-Zn-Ag VMS sulfide
with a $\sim$235 Ma Au-bearing pyrite provides two independently
well-constrained age brackets against which to test He retentivity and
closure behavior in pyrite across very different formation temperatures,
host lithologies, and post-depositional thermal histories.

\subsubsection{Kidd Creek, Ontario, Canada}

Kidd Creek is among the largest and most profitable VMS deposits in the
world, with past production and reserves exceeding 138.5 Mt at average
grades of 2.4 percent Cu, 6.5 percent Zn, 0.23 percent Pb, and 90 g/t Ag
\citep{hannington1999}. The deposit occurs within the Late Archean
Kidd-Munro assemblage of the Abitibi subprovince, Superior Province, hosted
at the top of a thickened pile of massive rhyolite flows, felsic
volcaniclastics, and coarse epiclastic rocks that attain a maximum
thickness of $\sim$300 m beneath the mine \citep{hannington1999}. The three
main stratiform orebodies (North, Central, and South) formed largely by
subseafloor infilling and replacement of permeable volcaniclastic rocks
within a subsiding, rift-related graben, rather than by high-temperature
venting at the seafloor \citep{hannington1999}. U-Pb zircon geochronology
brackets the ore-forming interval tightly: the oldest footwall rhyolite is
dated at 2716.1 $\pm$ 0.6 Ma, the onset of bimodal komatiite-rhyolite
volcanism at 2717 $\pm$ 1.1 Ma, and the youngest, uppermost quartz porphyry
unit at 2711.5 $\pm$ 1.2 Ma, indicating that mineralization took place
during at least one and possibly several million years of intermittent
volcanic activity \citep{hannington1999}. Oxygen isotope studies of the
host rhyolites show that, while primary igneous zircon retains magmatic
$\delta^{18}$O values (5.4\textperthousand{} $\pm$ 0.3\textperthousand{}), coexisting quartz
phenocrysts are pervasively shifted to much higher and more heterogeneous
values (14.2\textperthousand{} to 16.1\textperthousand{}), recording extensive subsolidus
recrystallization during the same long-lived hydrothermal convection that
deposited the sulfides \citep{king1997}.

The same 2.7 Ga host rocks that hosted this ancient hydrothermal system
now host the Kidd Creek Deep Fluid and Deep Life Observatory, 2.4 km below
surface, where fracture waters with mean residence times of 0.9-1.5 billion
years -- among the oldest groundwaters yet documented on Earth -- sustain a
low-biomass but metabolically active community of sulfate-reducing and
alkane-oxidizing microorganisms \citep{lollar2019}. The persistence of a
thermal regime compatible with life for roughly the past two billion years
in these fracture networks is a striking, if incidental, testament to how
long a deep, once-hydrothermally-active sulfide system can remain
hydrologically and chemically distinct from the surface -- a timescale of
isolation that dwarfs anything accessible to direct observation, but that
resonates with the closure and retentivity questions this chapter poses for
pyrite as a He geochronometer.

\subsubsection{Gympie, Queensland, Australia}

The Gympie pyrite comes from a very different setting: an accreted,
Early Permian bimodal (basalt-dacite) submarine volcanic sequence, the
Gympie Group, interpreted as a fragment of an intra-oceanic island arc that
rifted from the Gondwana margin before docking against the New England
Orogen of eastern Australia \citep{sivellwaterhouse1988}. Drillcore near
the Gympie Eldorado Gold Mine resolves at least nine sequentially emplaced
magmatic suites and a $>$2 km thick volcano-sedimentary succession, from the
basal tholeiitic Highbury Volcanics (Alma, Tozer, and Mary members, $>$700
m) upward through the more felsic andesite-dacite Rammutt Formation
($\sim$360 m), the alkaline Langton Dolerite, the South Curra Limestone, and
the carbonaceous, pyritic Tamaree Formation turbidites, capped by Early
Triassic Kin Kin beds \citep{sivellmcculloch2001}. This arc terrane is
inferred to have accreted to the Gondwana margin by about 220 Ma (Early to
Middle Triassic), based on cross-cutting granitoid relationships with
adjacent New England Orogen terranes \citep{sivellmcculloch2001}.

Gold-bearing pyrite-quartz mineralization at Gympie -- the ``Gympie vein
set'' -- occurs in fibrous extensional veins developed within sheared,
carbonaceous black slates, syntectonic with the cleavage-forming
deformation associated with the regional Hunter-Bowen orogeny
\citep{holcombe1997}. Alteration associated with this vein set has yielded
a K/Ar age of about 235 Ma, indistinguishable within the reported
$\sim$245-235 Ma range of K/Ar alteration ages obtained from the other
major gold deposits of the northern New England Fold Belt, including
Cracow \citep{holcombe1997}. Refractory dolerite dykes and later
microdiorite-shoshonite dykes cross-cut the mineralized structures and
bracket the timing of the auriferous quartz veining from below and above,
respectively, giving an independent field constraint on the age of
mineralization that is broadly consistent with the K/Ar alteration ages
\citep{sivellmcculloch2001}.

\section{Experimental Methods}\label{sec:ch3-experimental-methods}

\subsection{U-Th/He Geochronology}\label{sec:ch3-experimental-methods-geochronology}

\subsubsection{Sample Preparation}

Intact crystals of pyrite and arsenopyrite were extracted from samples from the Kidd Creek and Gympie deposits with a 
hammer and chisel at Caltech and SELFRAG high-voltage pulse fragmentation at the University of Queensland, respectively. 
Individual crystals were washed in a sequence of acetone, ethanol, and MilliQ water baths, then promptly dried using clean, lint-free wipes.
These samples, along with the hematite MI-43 standard from Gogebic, Michigan \citep{farleymckeon2015} and the Apatite standard from 
Durango, Mexico \citep{farley2000} were weighed, then encapsulated in silica glass ampoules similarly to the sulfides in section \ref{sec:ch3-experimental-methods-step-heats}, 
but without any additionaly buffer materials. An empty tube was also encapsulated to serve as a procedural blank. Photos and masses of the
prepared ampoules were recorded so that the helium and trace metal analyses could be correctly matched.

\subsubsection{Helium Analysis}

Encapsulated samples were wrapped in Ni-foil and loaded into a custom-built autosampler with the procedural blank loaded first. The autosampler was wrapped in heating tape, and was made
to warm the samples to 80-100°C for at least 12 hours in ultra-high vacuum before the first step heat to allow for atmospheric helium to escape. Initial 
efforts in method development showed that atmospheric helium contamination was significant (\todo{APATITE GAS FIGURE}). This atmospheric contamination is of the order of 10s of nccSTP helium per ampoule \citep{thyer2026}.
and is difficult to mathmatically correct for even when 3He can be accurately measured, because helium diffusion and permeation through silica glass fractionates 3He/4He ratios \citep{shelby1971,jones1953}, and the 
ampoules were each geometrically unique and therefore had different permeation rates. The procedural blanks showed that atmospheric helium contamination was reduced by a factor of 5-10, which was sufficient for these helium-rich samples with 100s of nccSTP helium (\todo{APATITE GAS FIGURE}). Future work may develop a more rigorous
approach to experimentally remove atmospheric helium contamination or to mathematically correct for it using the first order framework of \citet{thyer2026}.

After 12 hours, unheated furnace blank measurements were periodically taken until the measured value was similar to such blanks 
taken for batches of samples not encapsulated in silica glass, then samples were individually dropped into a resistance furnace and baked at 1000-1200°C for 15 minutes to degas their helium contents. 
Although pyrite and arsenopyrite are unusually retentive of helium at lower pressures, these temperatures were far above these phases' stability field, and so they were readily degassed. 
Helium abundances and isotope ratios were measured on a Thermo Fisher Helix SFT mass spectrometer. Gasses evolved during the bake were passed through a charcoal filter held at 77°K, then exposed to an SAES-10 getter and a heated 
titanium getter before being cryofocused at 14°K, then released at 34°K. Sensitivities were quantified using a calibrated 3He-enriched air standard several times throughout the analysis.

\subsubsection{Trace Metal Analysis}
The encapsulated, now degassed samples were removed from the furnace and were matched to their pre-bake photos and masses. Around 20\% of samples would be typically lost during the helium analysis process,
with ampoules either broken by the autosampler or by thermal stresses during the bakes. Samples bearing residual nickel-foil that had fused to the surface of the glass were scratched clean with a dental pick,
and the remaining glass ampoule was sonicated in a sequence of ethanol then MilliQ water baths. Samples were then placed in an acid-washed Teflon beaker and were broken underwater by an agate pestle, to prevent 
loss of sample material. The pestle was rinsed with MilliQ water between samples to avoid cross-contamination. Excess water was pipetted away from each sample and disposed of, then the remaining water was evaporated on a hot plate.
The vast majority of the samples were then spiked with a U-232 and Th-230, and vanadium enriched spike and digested in concentrated HNO\textsubscript{3} in closed Teflon beakers, leaving the glass intact. 
The uranium and thorium concentrations of the glass ampoules were measured to be $\approx$ 50 and 30 ppb by mass respectively, but this contaminant payload was only delievered when the glass was dissolved in HF; 
the uranium and thorium contents of concentrated nitric acid in extended contact with broken glass were below detection limits. 

After this point, the Kidd Creek pyrite and Durango Apatite samples were pipetted and diluted to an ICPMS-ready solution. The gympie samples, however, required further processing to dissolve fine U-Th bearing rutile inclusions. The gympie and MS-43 hematite samples, now mostly digested
in concentrated HNO\textsubscript{3}, were dried down and put back up in small volumes of concentrated HCl, then were loaded into Parr vessels and baked at 200°C for 24 hours. 
The Gympie samples were then dried down again and put back up in concentrated HNO\textsubscript{3} before the concentrated acid was transferred and diluted to an ICPMS-ready solution. U and Th contents were quantified with isotope dilution using an Agilent 8800 ICP-MS. 
Other trace elements were measured as ratios with spike vanadium. The U-Th contents of the spikes and the sensitivities of the other quantified trace metals relative to vanadium were continuously montiored by analyzing a standard solution of known composition that was spiked with the same U-Th-V spike as the samples.
The same empty tube procedural blank used for the helium analysis was also processed through the same acid digestion and ICPMS analysis, and the measured U, Th, and other trace metal contents were subtracted from the measured values of the samples. The U-Th contents of the procedural blank were below detection limits, and the other trace metals were typically 1-5\% of the measured values of the samples.


\subsection{Encapsulated Sulfide Step Heat
Experiments}\label{sec:ch3-experimental-methods-step-heats}

\textbf{Samples}

Nine helium step heat experiments of sulfides (pyrite, arsenopyrite,
chalcopyrite, galena, chalcocite, and sphalerite) sealed in silica glass
ampoules performed at the Noble Gas Laboratory at Caltech are discussed
here. Table 1 shows the dimensions, geometric factors, and encapsulated
materials of this work. All sulfides of this study were sourced from the
Caltech Mineral Collection from visually pure, well-formed euhedral
crystals with half prism widths greater than 200 microns, often larger
than a centimeter. Samples were broken into smaller pieces under ethanol
to avoid irreversible trapping of atmosphere \citep{cox2022}. Except
for the Rustenburg pyrite, all step heats measured 3He evolution from an
initially homogeneous distribution of spallogenic 3He. To produce a
homogeneous distribution of 3He, the samples were irradiated with 220
MeV protons for 24 hours at the Francis H. Burr Proton Therapy Center,
Massachusetts General Hospital, following standard procedures \citep{shuster2005}\textbf{.}

\begin{longtable}[]{@{}
  >{\raggedright\arraybackslash}p{(\linewidth - 12\tabcolsep) * \real{0.1951}}
  >{\raggedright\arraybackslash}p{(\linewidth - 12\tabcolsep) * \real{0.1431}}
  >{\raggedright\arraybackslash}p{(\linewidth - 12\tabcolsep) * \real{0.1198}}
  >{\raggedright\arraybackslash}p{(\linewidth - 12\tabcolsep) * \real{0.1219}}
  >{\raggedright\arraybackslash}p{(\linewidth - 12\tabcolsep) * \real{0.0741}}
  >{\raggedright\arraybackslash}p{(\linewidth - 12\tabcolsep) * \real{0.1942}}
  >{\raggedright\arraybackslash}p{(\linewidth - 12\tabcolsep) * \real{0.1123}}@{}}
\caption{Table 1: Silica Glass encapsulated sulfides that were the
subject of helium step heat experiments in this work. The Rustenburg
Pyrite experiment had evolved radiogenic 4He quantified by isotope
dilution using a Pfeiffer Prisma QMS 200 quadrupole mass spectrometer,
whereas evolved spallogenic 3He was quantified using a
Helix-MC-Plus\textsuperscript{10K} multi-collector mass spectrometer for
the remainder of the samples, which were proton-irradiated. Experiments
Spal230707 and CCT230707 contained no steps for which the uncertainty
due to helium permeation through the silica glass capsule could be
estimated, and therefore are discussed in a limited, qualitative
capacity only.}\tabularnewline
\toprule\noalign{}
\begin{minipage}[b]{\linewidth}\raggedright
Experiment Name and Phase
\end{minipage} & \begin{minipage}[b]{\linewidth}\raggedright
Sample Locality
\end{minipage} & \begin{minipage}[b]{\linewidth}\raggedright
Tube Inner x Outer Diameters
\end{minipage} & \begin{minipage}[b]{\linewidth}\raggedright
Geometric factor \emph{G} (\emph{eqn 2.2, Ch.1)}
\end{minipage} & \begin{minipage}[b]{\linewidth}\raggedright
VR (\emph{eqn 5.4, Ch.1)}
\end{minipage} & \begin{minipage}[b]{\linewidth}\raggedright
Encapsulated ``Buffer'' Materials
\end{minipage} & \begin{minipage}[b]{\linewidth}\raggedright
Mineral Half-prism Width (a)
\end{minipage} \\
\midrule\noalign{}
\endfirsthead
\toprule\noalign{}
\begin{minipage}[b]{\linewidth}\raggedright
Experiment Name and Phase
\end{minipage} & \begin{minipage}[b]{\linewidth}\raggedright
Sample Locality
\end{minipage} & \begin{minipage}[b]{\linewidth}\raggedright
Tube Inner x Outer Diameters
\end{minipage} & \begin{minipage}[b]{\linewidth}\raggedright
Geometric factor \emph{G} (\emph{eqn 2.2, Ch.1)}
\end{minipage} & \begin{minipage}[b]{\linewidth}\raggedright
VR (\emph{eqn 5.4, Ch.1)}
\end{minipage} & \begin{minipage}[b]{\linewidth}\raggedright
Encapsulated ``Buffer'' Materials
\end{minipage} & \begin{minipage}[b]{\linewidth}\raggedright
Mineral Half-prism Width (a)
\end{minipage} \\
\midrule\noalign{}
\endhead
\bottomrule\noalign{}
\endlastfoot
LVCOPy230630

Pyrite & Leadville, Colorado & \emph{Not Measured} & \emph{3000
cm\textsuperscript{-2}} & - & \textasciitilde0.2 mg Fused Sulfur &
\textasciitilde200 µm \\
BigLVCOPy230707 Pyrite & Leadville, Colorado & \emph{\hl{Not Measured}}
& \emph{\hl{3000 cm\textsuperscript{-2}}} & \hl{-} & \textasciitilde0.2
mg Fused Sulfur & \hl{\textasciitilde400 µm} \\
112123LVCOPyIrr Pyrite & Leadville, Colorado & 830 x 1550 µm & 1859
cm\textsuperscript{-2} & 2.5 & \textasciitilde{} 0.2 mg Fused Sulfur
\emph{(and \textgreater5 mg silica glass shards to separate sulfur melt
from pyrite)} & \textasciitilde300 µm \\
Rustenburg Pyrite

Pyrite (4He on Quad) & Rustenburg, South Africa & 520 x 810 µm & 6675
cm\textsuperscript{-2} & 1.4 & \textasciitilde0.2 mg Fused Sulfur & 106
µm \\
112123SalAspIrr

Arsenopyrite & Mitterberg, Austria & 530 x 970 µm & 4711
cm\textsuperscript{-2} & 2.3 & \textasciitilde0.5 mg Crystalline Realgar
(AsS) & \textasciitilde175 µm \\
GalTS230630

Galena & Tri-State District, Missouri & 620 x 920 µm & 5273
cm\textsuperscript{-2} & 1.2 & \textasciitilde0.2 mg Fused Sulfur &
\hl{\textasciitilde100 µm} \\
CCpTS230630

Chalcopyrite & Baxter Springs, Kansas & 540 x 1010 µm & 4382
cm\textsuperscript{-2} & 2.5 & \textasciitilde0.2 mg Fused Sulfur &
\textasciitilde200 µm \\
Spal230707*

Sphalerite & Chitwood, Missouri & 1110 x 1490 µm & 2205
cm\textsuperscript{-2} & 0.8 & \textasciitilde0.2 mg Fused Sulfur &
\textasciitilde200 µm \\
CCT230707*

Chalcocite & Messina, South Africa & 1100 x 1520 µm & 2044
cm\textsuperscript{-2} & 0.9 & \textasciitilde0.2 mg Fused Sulfur &
\textasciitilde200 µm \\
\end{longtable}

\textbf{Encapsulation}

Before encapsulating the sulfide samples, glass tubes were cleaned and
pulled into pipettes. 10 cm long segments of pure silica glass tubes
with an outer diameter of 4 mm and an inner diameter of 3 mm were
submerged and sonicated within acetone, ethanol, then MilliQ water baths
in sequence, and then were dried overnight in a 100°C oven. The glass
tubes were hand-pulled into thinner pipettes using a natural gas and
oxygen welding torch and the thin ends were fused closed.

Aliquots of sulfides were weighed and loaded into room temperature glass
tubes along with fused sulfur in the presence of an ionizing fan to
mitigate static. The arsenopyrite ampoule was loaded with crystalline
realgar (arsenic sulfide) instead of fused sulfur, and one pyrite
ampoule was loaded with silica glass pieces above and below the pyrite
aliquot to separate the pyrite from the sulfur liquid during the
experiment. To seal, the top of each pipette was interfaced with a
vacuum line and the tube was pumped to a pressure of \textless60
millitorr, which greatly aided the ease and quality of the final seal.
After attaining rough vacuum, a fine-flamed welding torch was used to
seal the top of each ampoule. To determine the geometric factors for
permeation of the ampoules, the inner and outer diameters of most
ampoules were measured using the measurement function of the Amscope
stereo microscope software, which was calibrated and tested during early
development of this method. \hl{Two experiments (LVCOpy230630 and
BigLVCOPy230707) were completed without these geometric measurements.
For these experiments, a geometric factor of 3000 cm\textsuperscript{-2}
was arbitrarily assigned to estimate uncertainties of the measurements
due to the convolution of the process of helium permeation through
silica glass. This assumption is unlikely to affect the interpretation
of the results}.

The sealed ampoules were placed in copper foil packets and mounted to a
K-type thermocouple within an ultra-high vacuum projector lamp heating
cell described in detail by \citet{farley1999}. Before the experiment
began, each ampoule was washed with ethanol, air-dried, loaded in the
lamp cell, pumped to high vacuum, and was held at 100°C for at least 30
minutes under dynamic vacuum to desorb ethanol, water, and other
atmospheric species on the ampoule, copper packet, and lamp cell.

\textbf{Helium Mass Spectrometry}

Heating steps and 3He measurements for all samples were automated using
custom LabView routines on a Helix-MC-Plus\textsuperscript{10K}
multi-collector mass spectrometer, or a Pfeiffer Prisma QMS 200
quadrupole mass spectrometer for the Rustenburg Pyrite ampoule. For all
ampoules, a Watlow temperature controller held the ampoule at the
temperature set point for a pre-programmed duration, after which the
ampoule was allowed to cool while the evolved gas was prepared for
measurement for about 25 minutes.

For the irradiated samples, Helium was purified by exposing evolved
gasses to a charcoal trap held in liquid nitrogen, a titanium sponge
getter, then a SAES NP10 getter before freely expanding into the
Helix-MC-Plus\textsuperscript{10K}. 3He was measured with a
pulse-counting electron multiplier, while 4He was detected on a Faraday
detector. the Helix-MC-Plus\textsuperscript{10K} mass spectrometer was
operating in peak hopping mode and the majority of the analysis time was
spent measuring 3He. Every few hours, the background was mass-scanned to
set the 3He and 4He measure points; each was taken to be a constant mass
offset from the background HD\textsuperscript{+} peak. These mass
offsets are precisely measured when tuning the
Helix-MC-Plus\textsuperscript{10K} for helium isotope analysis. Months
of helium measurements with the Helix-MC-Plus\textsuperscript{10K} at
Caltech has shown this approach to be reliable, reproducible, and
accurate for helium isotope analysis. Additionally, an electron
multiplier monitored atomic mass 3.5 to detect 4He overspray. 4He
signals were near-blank for all experiments of irradiated samples, and
no attempt is made here to interpret those data. 4He overspray likewise
did not affect 3He signals. Measurements of high 3He MM gas standard
were analyzed every 8 hours to track instrumental drift of 3He and 4He
sensitivities and to quantify internal reproducibility. These standards
show 3He was generally internally reproducible to within 1\% (2$\sigma$) for
rolling sets of 5 consecutive standards, except for the latest
measurements on the 112123SalAspIrr Arsenopyrite, during which the ion
source began to fail. Those measurements are excluded from further
analysis, and any helium evolved during those steps is likely to be a
very small proportion of the total helium contents of the sample
(\textless0.2\%) and thus would not significantly affect the calculation
of the total helium content of the sample.

Radiogenic 4He evolved during step heats of the encapsulated Rustenburg
Pyrite was quantified using isotope dilution with a high purity 3He
spike on a Pfeiffer Prisma QMS 200 quadrupole mass spectrometer. A
standard gas with a known amount of 4He was measured before and after
the experiment, to monitor drift in the relative sensitivities of 3He
and 4He and to estimate reproducibility, which was about 6\% (2$\sigma$). After
each step, evolved gas from the ampoule in the lamp heating cell was
allowed to mix with the 3He spike and the mixture was frozen on charcoal
held at 14°K within a cryostat, then helium was selectively released to
the mass spectrometer by raising the temperature to 36°K.

\textbf{Scanning Electron Microscope Study}

The methodology for the SEM study of the sulfides is described in detail in \ref{sem-imaging-of-buffered-and-decomposed-sulfides}.

\subsection{Analytical Methods For Step Heat Experiments}\label{sec:ch3-analytical-methods}

Similarly to \citet{yakubovich2019}, these helium step heat experiments were designed according to the 
incorrect assumption that the silica glass ampoule was
effectively ``helium transparent'', or in other words that helium
permeation through the silica glass ampoule was nearly instantaneous
with respect to the rate of helium evolution from the encapsulated
sulfides. Thus, permeation steps (as discussed in chapter 1 and 2), were
rarely used. The step yields for steps at high temperatures were approximated
as the the measured value, with permeation lag constrained by upper bounds following the 
"Treatment as Uncertainty" approach of \citet{thyer2026}. This approach
requires no permeation steps or knowledge of the helium content within ampoule from previous
steps, instead calculating the maximum helium content remaining within
the ampoule for a given measurement, provided the step was long and hot
enough for the tube to have reached a steady state concentration and
assuming that the rate of helium evolution from the encapsulated mineral
is constant. The steady state assumption is taken to be effectively
valid once the ampoule loses memory of the helium content remaining from
the previous step. According to eqn.~\eqref{eq:s6-8} in chapter 1, only 1.8\% of the
difference between the previous step's remaining helium content and the
current step's steady state helium content remains when
\(\lambda t = 4\). In this work, the steady state assumption is only
applied for steps long and hot enough for \(\lambda t \geq 4,\ \)where
\(\lambda = KGRT\), \emph{K} is the temperature-dependent permeation
coefficient for 4He or 3He through silica glass, \emph{G} is the ampoule
geometric factor, \emph{R} is the gas constant, \emph{T} is the
temperature in Kelvin, and \emph{t} is the time in seconds. The 4He
permeation coefficient \emph{K} is taken from \citet{swets1961}; and
these coefficients are multiplied by the permeation coefficient ratio
from \citet{jones1953} to yield the 3He permeation coefficient.

Many heating steps, especially near the beginning of the experiments
where lower temperature steps were used, are not interpreted or plotted
in this work because \(\lambda t < 4\); however, some steps with
\(\lambda t < 4\) are used when eqn.~\eqref{eq:steadystate} is used to place upper bound
constraints on preceding steps, so that eqn.~\eqref{eq:s6-8} can in turn be used to
quantify the uncertainty on the yields step that has not reached steady
state. 

This method of calculating uncertainties produces an asymmetric upper
and lower uncertainty for each step. These assymetric uncertainties are propagated 
through the diffusivity inversion of \citet{fechtigkalbitzer1966} with eqn.~\eqref{eq:diffusivity-bounds}
in Chapter 2.

In general, the 2$\sigma$ analytical uncertainties as calculated using the
reproducibility of standards is greatly exceeded by the uncertainties
due to the permeation lag of the ampoule as calculated above. For the
Rustenburg pyrite experiment and LVCOpy230630, the uncertainty due to
the permeation lag is less than a factor of 10 greater than the
analytical uncertainty, and so the 2$\sigma$ analytical uncertainty for a given
measurement is added to the uncertainty by adding it to the upper bound
helium content of the tube, \(n_{t\ \ }^{UB}\), for that step. For the
remainder of the experiments, the analytical uncertainty is less than
the uncertainty from the permeation lag by a factor of 10-100, and so
the analytical uncertainty is neglected. No uncertainty in the total
yields in incorporated.

Two of the experiments contained no steps for which
\(\lambda t \geq 4\): Spal230707 and CCT230707 for sphalerite (ZnS) and
chalcocite (Cu\textsubscript{2}S) respectively. No step heat results are
reported here for those experiments, but they are discussed
qualitatively.

\section{Results}\label{sec:ch3-results}

\subsection{Effectiveness of the Ampoule Method in Stabilizing Sulfides}\label{ch3-sem-imaging-tube-buffered-sulfides}
The methodology for the heating treatments and SEM imagery of pyrite and arsenopyrite 
surfaces following thermal treatments in different conditions is discussed in \ref{sem-imaging-of-buffered-and-decomposed-sulfides}.
The SEM imagery reveals deformed, variably raised, \textasciitilde10 µm
cracked surfaces on pyrite and marcasite that are heated in open
\hl{air, and for arsenopyrite heated in a closed capsule without a
realgar buffer (Figures 3.2\ldots)}. Pyrite and marcasite samples that
are heated in ultra-high vacuum developed recessive \textasciitilde1 µm
rounded craters. Contrarily, images of thermally treated, encapsulated
pyrite and arsenopyrite with added fused sulfur or realgar show no
obvious distinction from the control group: both of which showcase
conchoidal fracture \hl{(Figures 3.2\ldots)}. The Rustenburg pyrite
sample represents a potential exception; its pervasive fractures appear
to have partially annealed after thermal treatment at high sulfur
fugacity \hl{(Figures 3.2\ldots)}. The buffered, encapsulated marcasite
sample shows a raised layer of $\leq$2 µm angular (presumably pyrite)
crystals growing atop its surface; evidence of phase transformation is
present in all three heat-treated marcasite populations, demonstrating
that the encapsulation method is not effective in preventing phase
breakdown at elevated temperatures for marcasite. While these
experiments are not conclusive evidence that thermal phase
transformation is prevented for pyrite or arsenopyrite, they consistent with that 
interpretation, encouraging its application for helium step heat experiments.

\subsection{Sulfide Step Heat Experiments}\label{ch3-sulfide-step-heat-experiments}

All apparent diffusivities calculated from the helium step heat data are
plotted together in figure 3.3 along with a reference line for helium
diffusion within Durango apatite with a half-prism width of 150µm
\citep{farley2000}. Figure 3.3 qualitatively shows remarkably high helium
retentivity at high temperatures for pyrite and arsenopyrite. For no
pyrite or arsenopyrite helium step heat experiment did the final
cumulative fractional helium yield exceed 7\%.

The irradiated Leadville pyrite was the most retentive sample of the
sulfide suite. Three separate helium step heat experiments were
conducted on this sample, two with varying grain size and a third that
attempted to isolate the pyrite from the sulfur liquid by packing glass
shards above and below the pyrite grains within the ampoule \hl{(figs.
3.0xx)}. The Arrhenius plots for the two experiments with varying size
fractions and no attempt to isolate the pyrite from the liquid sulfur
were similar to one another but distinct from the third in which the
pyrite was isolated from the sulfur, which was the experiment that
showed the highest helium retentivity of the sulfide sample suite
(figure 3.3; figure 3.4a). The three experiments shared a similar slope,
and a regression was fit to a linear region of 9 measurements of the
LVCOpy230630 experiment with an R\textsuperscript{2} of 0.998, yielding
an apparent activation energy of 47.6 Kcal/mol (figure 3.4a). The delta
plot constructed with this activation energy showed a monotonic decrease
in delta values for the experiments that did not attempt to separate the
pyrite from the sulfur liquid. The more retentive ampoule that held the
pyrite in the middle of the tube away from the majority of the sulfur
liquid also preserved this general trend, but showed a slight increase
in delta values midway through the experiment (figure 3.4b).

The octahedral Rustenburg pyrite is less retentive of helium than the
pyritohedral Leadville pyrite. On the down-temperature retrogrades of
heating cycles, the Rustenburg pyrite shows clearly developed
reproducible lines with the unusually high implied activation energy of
71.8 Kcal/mol (figure 3.5a: R\textsuperscript{2} of 0.998 on 5
measurements). This feature is completely absent from the three
experiments of proton-irradiated Leadville pyrite (figure 3.4a). The
delta plot constructed using this activation energy shows a monotonic
decrease in delta values throughout the duration of the experiment
(figure 3.5b).

The encapsulated arsenopyrite helium step heat experiment produced an
apparent diffusivity Arrhenius plot that showed generally high helium
retentivity, but less so than the Leadville pyrites (figure 3.3). There
were no representative linear regions produced by this experiment that
were suitable for linear regression, although the general slope was
similar to the Leadville pyrites' and so it was also used as a reference
for this experiment (figure 3.6a). The delta plot using this reference
showed complex behavior, where high delta values (far beyond the
uncertainty related to the permeation lag of the ampoule itself) were
observed for low temperature steps following high temperature steps
(figure 3.6b).

The majority of the helium total yield was evolved during the step heat
experiment of the encapsulated galena, in stark contrast to the
comparatively retentive pyrites and arsenopyrite. The reference line
from the chalcopyrite experiment showed a similar activation energy to
the activation energy implied by curved retrogrades of the galena
experiment (figure 3.7a), so it was used as the reference for the delta
plot of galena. This delta plot shows a monotonic decrease in delta
values as the experiment proceeds (figure 3.7b).

Step heat data step heat experiments for chalcopyrite (CCpTS230630) and
chalcocite (CCT230707) produced macroscopic evidence of phase
transformation, including the development of an iridescent blue sulfide
or oxide layer on both specimens and dramatic crystal growth on the
ampoule walls for the chalcocite experiment \hl{(figs. 3.0xx).} Although
permeation lag was not sufficiently constrained to calculate apparent
diffusivities for the chalcocite experiment, the data for CCT230707 are
consistent with a primary phase-breakdown control of helium migration at
temperatures $\leq$300°C.

A linear regression is fit to 9 measurements defining a linear region of
the Arrhenius plot of chalcopyrite sample CCpTS230630 for use as a
reference for the delta plot (figure 3.8a). This regression has an
R\textsuperscript{2} value of 0.9811, an implied activation energy of
41.8 kcal, and an implied closure temperature of 128°C assuming a
cooling rate of 10 °C/Ma. The delta plot generated using this reference
line for chalcopyrite clearly shows a dramatic mid-experiment increase
in diffusivity following steps above 400°C, followed by a continuous
decrease in delta values (figure 3.8b), similarly to that observed for
the dehydrating goethite sample Win-06-01b \citep{farley2024}.

\section{Discussion}\label{sec:ch3-discussion}

\subsection{Remarkable Helium Retentivity of Pyrite}\label{ch3-discussion-retentivity}


\subsection{Inferring Bulk Helium Migration Behavior from Low Fractional
Yields}\label{ch3-low-fractional-yields}

Many of the step heats on encapsulated sulfides presented here show very
high helium retentivity, and for those experiments only a very small
fraction of the total yield evolved (\textless5\% for the arsenopyrite
and Rustenburg pyrite; \textless2\% for the Leadville pyrite), similarly
to previous work on highly helium retentive phases such as baddeleyite
\citep{baughmanflowers2018}. For these sulfides, a
mechanism responsible for bulk helium loss besides desulfurization
cannot be identified. For some well-studied phases like zircon and
titanite, Arrhenius plots of the earliest fractions of evolved helium
often reflect processes or conditions that are accessory to the bulk
mechanism for helium migration through the mineral, such as radiation
damage annealing or surface roughness \citep{reinersfarley1999,
guenthner2013}, where the primary mechanism for helium
migration through the pristine mineral is volume diffusion. These
apparent departures from the bulk helium migration behavior can be
caused by properties of natural samples, like surface roughness or
fine-scale cracks and fractures \citep{reinersfarley1999} that produce
complex grain size distributions even in carefully sieved, (apparently)
single grains. Contrastingly, some other complexities in the helium loss
behavior at early loss fractions can be ephemeral and altered during
experimental heating, such as the annealing of radiation damage \citep{shusterfarley2009}. For radiation damaged samples, understanding the
migration behavior of helium before and during annealing is essential to
understand the sample's retentively in nature (\hl{citations}) and
therefore the earliest fraction of helium evolved in step heats is the
focus of many studies (\hl{citations}).

The fraction of total helium evolved before the predominant migration
mechanism's behavior can be characterized in helium step heats varies:
\citet{guenthner2013} show that significantly non-zero delta or
ln(r/r0) values persist in zircon until about 2\% of the total helium is
degassed for typical natural samples; \citet{reinersfarley1999} observe
apparent diffusivities of helium in titanite converging to linearity in
titanite after 5\% of the total helium is evolved. Similar experiments
on proton-irradiated synthetic samples seem not to be marred by the same
complexities: helium step heat data for synthetic monazite converge
nearly immediately to a linear array at very loss cumulative helium loss
fractions \citep{farley2007}.

The step heat experiments on pyrite and arsenopyrite in this study would
have benefitted from longer steps at high temperature to degas more
helium to better understand the bulk behavior of helium, but even at
small loss fractions some key observations about helium loss in sulfides
can be made.

\subsection{Comparison of Helium Step Heat Experiments of Natural Pyrites}\label{ch3-compare-pyrites-discussion}

In this work, four encapsulated, sulfur-buffered aliquots of pyrite
taken from two localities are subjected to helium step heat experiments.
Step heats for the three irradiated pyrites from Leavldville, CO are
broadly similar, but contrast markedly from step heats of non-irradiated
pyrite from Rustenburg, South Africa, therefore differences between
pyrites of different localities is discussed before the variation within
the aliquots from Leadville, CO.

Radiogenic helium evolution in step heats of octahedral pyrite from
Rustenburg, South Africa show that this sample is less retentive than
the pyritohedral pyrite from Leadville, Colorado \hl{(fig citation}).
Apparent diffusivities of these samples also have fundamentally distinct
geometries when plotted in an arrhenius plot. The Rustenburg pyrite
shows a striking feature of very steep slopes (corresponding to an
activation en\hl{ergy of \_\_} ) on retrograde-prograde heating cycles
at temperatures lower than the previous maximum temperature step heat,
whereas the maximum apparent diffusivity for a given temperature seemed
to follow a much shallower slope (lower activation \hl{energy of )},
which was broadly similar to the slope and activation energy of the
Leadville, CO pyrite. The geometry of the step heats in an arrhenius
plot for the irradiated Leadville, CO pyrite is generally simpler:
generally slowly decreasing in a delta plot, or for isothermal steps
(fig \hl{citation}) as would be expected for helium release governed by
MDD systematics (\hl{citation}).

Natural pyrite is known to develop fractures and crystalline
dislocations in metamorphic and experimental conditions \citep{gill1969,
mcclayellis1983,barrie2011}. Abundant fractures of
micron-scale widths were clearly visible in the Rustenburg pyrite of
this study (fig 3.0a). Those pyrites were surprisingly easy to crush,
presumably breaking easily along fine reexisting fractures. Although
pyrites from Leadville are known to be associated with faulted
dolomites, especially along deformation planes which likely served as
pathways for mineralizing hydrothermal fluids \citep{nash1979}, the samples
used here were exceptionally sturdy and visually homogeneous under SEM.
Rare veins and coatings of dolomite within and around the pyrites were
present but were avoided when preparing the sample for irradiation.

Within this work, a pipe/short circuit diffusion mechanism is favored to
explain the helium loss behavior of the Rustenburg pyrite, where helium
populates the sample's fine fractures at increasing temperatures due to
the primary helium migration mechanism for a system of encapsulated
pyrite and sulfur fluid: volume diffusion or fine-scaled
dissolution/reprecipitation of pyrite within sulfur fluid. This primary
mechanism is interpreted to be the same for all pyrite samples
regardless of locality. On retrograde heating steps, short circuit
diffusion of helium out of fine fractures in the Rustenburg pyrite is
hypothesized to become the rate-limiting process.

\subsection{Leadville, CO Pyrite: Evidence for a Sulfur Liquid Control for
Helium Loss in Step Heat Experiments}{ch3-leadville-pyrite-sulfur-liquid-control}

Three helium step heat experiments for the irradiated Leadville, CO
pyrite are discussed here. Grain size was varied by a factor of 2
between a half-prism width of 200 microns for LVCOPy230630 and 400
microns for LVCOPy230707 to investigate the dependence of grain size on
helium migration in pyrite. The aim of the third experiment,
112123LVCOPy, was to test the effect of isolating the pyrite from the
liquid sulfur during the experiment to investigate the role of
dissolution and (re)crystallization of pyrite by the ampoule's sulfur
fluid. Silica glass shards were added to this ampoule below and above
the pyrite aliquot, forcing the pyrite grains to remain in the middle of
the ampoule, and not at the edges/bottom, where sulfur liquid is
presumed to accumulate. Despite the similarities of the size fraction of
the glass-filled 112123LVCOPy (300 micron half-prism width) to the other
two aliquots of pyrite, it yielded apparent diffusities that were
\hl{markedly more retentive}. This suggests that reaction between the
sulfur liquid and pyrite was the primary control for helium release from
the Leadville samples. This effect is likely to be important for helium
release in the Rustenburg pyrite as well.

Although iron solubility in sulfur liquids has only been sparsely
studied, ferrous iron is known to be only modestly soluble in sulfur
liquids. \citet{clark1960} reported that sulfur liquid can be up to 0.12\%
iron. Later, \citet{barton1969} reproduced this measurement and reported iron
contents in pyrite-arsenopyrite saturated S-As liquids ``less than 0.5,
and probably less than 0.2, wt.\%''. In the step heat experiments of
this work, the mass of the added sulfur was not recorded but was always
small: less than or equal to the mass of the pyrite aliquot. By simple
mass balance, only a maximum of 1.5\% of the pyrite would be dissolved
in an equal mass of sulfur fluid, however, if iron solubility in sulfur
liquid is highly temperature dependent within the range of temperatures
used in these experiments (24-500°C), then this effect could be
multiplicative with each thermal cycle: the sulfur liquid could
cyclically dissolve and recrystallize pyrite during thermal cycling. How
this dissolution-recrystallization cycle might affect step heat results
is unclear, but intuitively it would likely libertate helium near the
grain surface early in the experiment, producing artificially high
apparent diffusivities for 3He cumulative fractional yields less than
1.5\%. This effect could potentially have been a primary control on the
degassing behavior of all the irradiated pyrite within this work, but it
is unlikely to be significant for the majority of the 4He step heat data
of the Rustenberg pyrite, which proceeded to beyond 4\% total cumulative
fractional yield.

The apparent diffusivities calculated for the two helium step heat
experiments of Leadville pyrite with varying half-prism width,
LVCOPy230630 and LVCOPy230707, both exhibit the same slope with
monotonically decreasing delta values (fig \hl{citation}). Intuitively,
the smaller aliquot contains higher delta values at early helium
fractional yields. These observations can be explained with MDD
systematics \citep[e.g.][]{farley2018} but they are also consistent with a
fluid moderated dissolution-recrystallization mechanism discussed above,
which would also proceed faster for smaller grains which expose more
surface area for the same mass of material. Similarly to volume
diffusion, for which irregular surfaces are degassed first \citep{reinersfarley1999}, dissolution-recrystallization would affect irregular
surfaces first, leading to the intuition that smaller grains will
contain more helium in less-retentive domains (at the grains' surfaces)
than larger grains with less surface area for the same mass.

\textbf{Relevance of Helium Step Heat Data of Encapsulated,
Sulfur-Buffered Sulfides in Nature}

The intention of these sulfur-buffered helium step heats was to
stabilize sulfides against thermal decomposition in the harsh vacuum
conditions typical in helium step heat experiments \citep{farley2024,
farley1999}. The data reported here for pyrite suggest that
sulfur-fluid moderated dissolution and recrystallization are governing
helium loss from highly crystalline, irradiated encapsulated pyrite in
the laboratory. These data therefore are not applicable to thermally
activated helium migration from pyrite (which is most often assumed to
be the cause of helium loss in nature and indeed is the basis for
thermochronometry: \citealp{ketcham2005}; \citealp{farley2002}). These data are still
useful, as they represent upper bound limits for rates of thermal helium
loss by volume diffusion or whatever thermally activated helium loss
mechanism is most important for pyrite in nature. In order for the
relatively small effect of liquid-pyrite dissolution-recrystallization
to be the primary driver of thermal helium loss for pyrite in the
laboratory, it must be a faster process than thermal helium migration
through crystalline pyrite.

For encapsulated sulfide step heat experiments in which high fractions
of helium (\textgreater10\%) are evolved, a sulfur liquid mediated
dissolution-recrystallization mechanism for helium loss is implausible
if the studied phase is not significantly visually altered. For example,
all of the 3He of Spal230707 (sphalerite), and most 3He of GalTS230630
(galena) were evolved during the step heats, but the grains remained
visually unchanged. For these experiments, helium was clearly lost by
thermally-activated processes from stable sulfides.

Given the high helium retention of pyrite and arsenopyrite shown in this
work, the predominant helium loss mechanism for these common sulfides in
nature may indeed be recrystallization. For example, pyrite metamorphic
conditions commonly recrystallize and may either shrink or grow \citep{craig1998,mcclayellis1983}. Until more work has characterized helium
loss (or gain) for recrystallizing pyrite or arsenopyrite in metamorphic
conditions, helium should cautiously be considered to have behaved as an
open or partially open system in sulfides during metamorphic events.

\textbf{Thermal Helium Migration Systematics for Other Sulfides}

\emph{Chalcopyrite and Chalcocite}

Neither copper-bearing sulfide mineral studied here was successfully
stabilized by encapsulation with sulfur, so these results do not provide
insight into helium migration within copper sulfides in nature. Still,
they are unlikely to be promising targets for U-Th/He geochronology or
helium isotope studies as there exists many copper-bearing sulfide
phases stable at temperatures between 200-500°C, for which phase
stability and Cu/Fe ratios are sensitive functions of temperature,
composition, and sulfur fugacity \citep{zhao2014,chenharvey1975,yundkullerud1966} and therefore are highly likely to have
reacted during cooling, reheating, or metasomatic/diagenetic events
after crystallization.

\emph{Galena}

Results from encapsulated Galena subjected to helium step heat
experiments were consistent with a multiple diffusion domain helium loss
mechanism. As expected for an MDD system, the delta plot for galena
monotonically decreases as the fractional yield increases (fig
\hl{citation}). Although no closure temperature can be inferred from
this data (because the domain size is variable), it qualitatively is
more retentive than (\hl{grain size)} apatite, suggesting that galena
that crystallized on the surface and experienced little heating due to
burial may be suitable for U-Th/He geochronology, similar to how
goethite and hematite are suitable for dating surface processes despite
the complexities associated with helium migration systematics in MDD
systems.

\textbf{Rustenburg: Pipe/short circuit diffusion mode}

\textbf{(JKW Lee: K2/K1 correlated with temperature, retrogrades see
grain size increase by annealing. (}

\textbf{Arsenopyrite}

No helium step heat data for sulfides plotted in figure 3.3 provide
evidence for a simple single domain volume-diffusion model, except
perhaps for the lone translucent sulfide of this study: sphalerite. The
sphalerite experiment (Spal230707) may allow a simple single domain
volume diffusion model. Unfortunately, the heating schedule of this
study was insufficient to deconvolve the process of helium permeation
through fused silica from helium evolution of sphalerite

Helium step heat experiments for galena (GalTS230630) and pyrite
(LVCOPy230630, BigLVCOPy230707, 112123LVCOPyIrr, Rustenburg Pyrite) are
all consistent with a multi-diffusion domain

Arsenopyrite: Arsenic degrades symmetry, moves from monoclinic to
triclinic with more arsenic, increases unit cell volume. \citep{morimotoclark1961}

Arsenopyrite becomes more arsenic rich at higher temperatures \citep{barton1969,morimotoclark1961}

\section{Conclusions}\label{sec:ch3-conclusions}
\todo{Summarize He retentivity findings and geochronologic implications.}
